\section{A Simple Example of Goodharting}
\label{appendix:simple_goodharting_example}

In ~\Cref{sec:explain_goodhart}, we motivated our explanation of Goodhart's law using a simple MDP $\mathcal{M}_{2, 2}$, with 2 states and 2 actions, which is depicted in ~\Cref{fig:mdp_graph}. We assumed $\gamma=0.9$,  and uniform initial state distribution $\m$.

\begin{figure}[h]
    \centering
        \begin{tikzpicture}[shorten >=1pt, node distance=3cm, on grid, auto]
            % States
            \node[state] (S0) {$S_0$};
            \node[state, right=of S0] (S1) {$S_1$};

            % Edges
            % For state 0
            \path[->, orange] (S0) edge[loop left, looseness=8] node[left] {$\frac{9}{10}$} (S0); % a=0, s=0->s=0
            \path[->, orange] (S0) edge[bend left, pos=0.5] node[above] {$\frac{1}{10}$} (S1); % a=0, s=0->s=1
            \path[->, purple, dotted] (S0) edge[loop right, looseness=8] node[right] {$\frac{1}{10}$} (S0); % a=1, s=0->s=0
            \path[->, purple, dotted] (S0) edge[bend right=120, pos=0.5] node[below] {$\frac{9}{10}$} (S1); % a=1, s=0->s=1

            % For state 1
            \path[->, orange] (S1) edge[bend left, pos=0.5] node[below] {$\frac{5}{10}$} (S0); % a=0, s=1->s=0
            \path[->, orange] (S1) edge[loop left, looseness=8] node[left] {$\frac{5}{10}$} (S1); % a=0, s=1->s=1
            \path[->, purple, dotted] (S1) edge[bend right=120, pos=0.5] node[above] {$\frac{8}{10}$} (S0); % a=1, s=1->s=0
            \path[->, purple, dotted] (S1) edge[loop right, looseness=8] node[right] {$\frac{2}{10}$} (S1); % a=1, s=1->s=1
    \end{tikzpicture}
    \caption{MDP $\mathcal{M}_{2, 2}$ with 2 states and 2 actions. Edges corresponding to the action $a_0$ are orange and solid, and edges corresponding to $a_1$ are purple and dotted.}
    \label{fig:mdp_graph}
\end{figure}
We have sampled three rewards $\R_0, \R_1, \R_2: \SxA \to \RR$, implicitly assuming that $\R(s, a, s') = \R(s, a, s'')$ for all $s', s'' \in \S$.

\begin{figure}[h]
    \centering
    \begin{subfigure}[t]{0.25\textwidth}
    \centering
    \begin{tabular}{|c|>{\color{orange}}c|>{\color{purple}}c|}
    \hline
    $\R_0$ & $a_0$ & $a_1$ \\
    \hline
    $S_0$ & 0.170 & 0.228 \\
    \hline
    $S_1$ & 0.538 & 0.064 \\
    \hline
    \end{tabular}
    \caption{Reward 0}
    \end{subfigure}
    \hspace{0.1\textwidth}
    \begin{subfigure}[t]{0.3\textwidth}
    \centering
    \begin{tabular}{|c|>{\color{orange}}c|>{\color{purple}}c|}
    \hline
    $\R_1$ & $a_0$ & $a_1$ \\
    \hline
    $S_0$ & 0.248 & 0.196 \\
    \hline
    $S_1$ & 0.467 & 0.089 \\
    \hline
    \end{tabular}
    \caption{Reward 1}
    \end{subfigure}
    \hspace{0.1\textwidth}
    \begin{subfigure}[t]{0.3\textwidth}
    \centering
    \begin{tabular}{|c|>{\color{orange}}c|>{\color{purple}}c|}
    \hline
    $\R_2$ & $a_0$ & $a_1$ \\
    \hline
    $S_0$ & 0.325 & 0.165 \\
    \hline
    $S_1$ & 0.396 & 0.114 \\
    \hline
    \end{tabular}
    \caption{Reward 2}
    \end{subfigure}
\caption{Reward tables for $\R_0, \R_1,\R_2$}
\end{figure}

We used 30 equidistant optimisation pressures in $[0.01, 0.99]$ for numerical stability. The hyperparameter $\theta$ for the value iteration algorithm (used in the implementation of MCE) was set to $0.0001$.